\documentclass[aspectratio=169]{beamer}
\usetheme{metropolis}
\usepackage{minted}
\usepackage{booktabs}
\usepackage{tikz}
\usetikzlibrary{arrows.meta, positioning, shapes.geometric}

\setminted{fontsize=\footnotesize, breaklines=true}

\title{APPy: Annotated Parallelism for Python on GPUs}
\subtitle{Work in Progress}
\author{}
\date{}

% ---------------------------------------------------------------
\begin{document}
\maketitle

% ---------------------------------------------------------------
\section{Motivation}

\begin{frame}{Why GPU Acceleration Matters}
  Scientific Python workloads keep growing in scale and complexity:
  \begin{itemize}
    \item Simulation \& numerical computing
    \item Data analysis \& processing pipelines
    \item ML and non-ML scientific workloads
  \end{itemize}
  \vspace{1em}
  GPUs offer massive parallelism — but Python's default execution model is \textbf{sequential}.

  \vspace{1em}
  \begin{alertblock}{The core problem}
    Most Python programs \emph{could} benefit from GPU acceleration, but the path there is not easy.
  \end{alertblock}
\end{frame}

\begin{frame}{Existing Path 1: Library-Based Acceleration}
  \textbf{Examples:} CuPy, PyTorch ops, JAX ops, cuBLAS
  \vspace{0.5em}

  \begin{columns}
    \column{0.48\textwidth}
    \textbf{Pros}
    \begin{itemize}
      \item Easy to use
      \item Pre-optimized kernels
      \item Great for standard linear algebra \& tensor ops
    \end{itemize}

    \column{0.48\textwidth}
    \textbf{Cons}
    \begin{itemize}
      \item Many workloads cannot be expressed using predefined operators
      \item Irregular loops, custom control flow, stencil codes $\to$ not supported
      \item Must rewrite algorithms to fit the library's operator set
    \end{itemize}
  \end{columns}
\end{frame}

\begin{frame}{Existing Path 2: Custom GPU Kernels}
  \textbf{Examples:} CUDA, OpenCL, Numba CUDA
  \vspace{0.5em}

  \begin{columns}
    \column{0.48\textwidth}
    \textbf{Pros}
    \begin{itemize}
      \item Full control over execution
      \item Potential for highest performance
    \end{itemize}

    \column{0.48\textwidth}
    \textbf{Cons}
    \begin{itemize}
      \item Must rewrite Python logic in a GPU kernel language
      \item Manual memory management
      \item Debugging: divergence, race conditions, shared-memory issues
      \item Hard to maintain correctness + portability
    \end{itemize}
  \end{columns}
\end{frame}

\begin{frame}{Existing Path 3: Triton}
  Triton's idea: write GPU kernels in a Python-like DSL that compiles to optimized GPU code.
  \vspace{0.5em}

  \begin{columns}
    \column{0.48\textwidth}
    \textbf{Strengths}
    \begin{itemize}
      \item High performance for dense, matrix/tensor-style workloads
      \item Minimal boilerplate vs.\ CUDA
      \item Strong auto-tuning infrastructure
    \end{itemize}

    \column{0.48\textwidth}
    \textbf{Limitations}
    \begin{itemize}
      \item Requires learning Triton's DSL — not pure Python
      \item Still kernel-centric: users must manage blocks, program IDs, memory loads
      \item Less natural for task-parallel or irregular workloads
    \end{itemize}
  \end{columns}
\end{frame}

\begin{frame}{The Gap in the Landscape}
  \begin{center}
  \begin{tikzpicture}[node distance=1.2cm]
    \node[draw, rounded corners, fill=mLightBrown!30, minimum width=3cm, minimum height=0.8cm] (lib) {Libraries (CuPy, JAX)};
    \node[draw, rounded corners, fill=mLightBrown!30, minimum width=3cm, minimum height=0.8cm, right=1.5cm of lib] (triton) {Triton DSL};
    \node[draw, rounded corners, fill=mLightBrown!30, minimum width=3cm, minimum height=0.8cm, right=1.5cm of triton] (cuda) {CUDA / OpenCL};

    \node[below=0.5cm of lib] {\small Too restrictive};
    \node[below=0.5cm of triton] {\small Kernel-centric DSL};
    \node[below=0.5cm of cuda] {\small Too low-level};

    \node[draw, rounded corners, fill=mLightGreen!50, minimum width=4cm, minimum height=0.9cm,
          below=1.8cm of triton] (appy) {\textbf{APPy} ← this gap};
    \draw[->, thick, dashed] (lib.south) -- ++(0,-0.7) -| (appy.north west);
    \draw[->, thick, dashed] (cuda.south) -- ++(0,-0.7) -| (appy.north east);
  \end{tikzpicture}
  \end{center}

  \vspace{0.5em}
  No simple way to express GPU parallelism using \textbf{regular Python loops} with \textbf{minimal annotation},
  preserving control flow across \textbf{multiple GPU paradigms}.
\end{frame}

% ---------------------------------------------------------------
\section{APPy: Prior Work (CC 2024)}

\begin{frame}[fragile]{APPy v1: OpenMP-Style Annotations for Python GPUs}
  \textbf{Key idea:} Write regular sequential Python, add a single pragma, get a GPU kernel.

  \begin{minted}{python}
  #pragma parallel for
  for i in range(N):
      C[i] = A[i] + B[i]
  \end{minted}

  \vspace{0.5em}
  No CUDA, no OpenCL, no Triton DSL required.

  \vspace{0.5em}
  Worked well for:
  \begin{itemize}
    \item[$\checkmark$] Dense, vectorizable loops
    \item[$\checkmark$] Regular data-parallel computations
    \item[$\checkmark$] Scientific workloads matching standard patterns
  \end{itemize}
\end{frame}

\begin{frame}{Where APPy v1 Hit a Limit}
  The CC 2024 model assumes a \textbf{single parallelism model}: every iteration is independent and maps to one GPU thread.

  \vspace{1em}
  But many real workloads:
  \begin{itemize}
    \item Contain \emph{inner} collective operations (dot products, reductions over slices)
    \item Have iterations that benefit from \emph{thread cooperation} within a group
    \item Are expressed most naturally with nested parallel structure
  \end{itemize}

  \vspace{1em}
  \begin{alertblock}{The bottleneck}
    Forcing everything into independent-thread parallelism is itself an abstraction leak.
    Users must restructure their code to fit the model.
  \end{alertblock}
\end{frame}

% ---------------------------------------------------------------
\section{The Dual Programming Model}

\begin{frame}[fragile]{APPy v2: A Unified Dual Parallelism Model}
  \textbf{One annotation. Two execution modes. Compiler decides.}

  \vspace{1em}
  \begin{columns}
    \column{0.48\textwidth}
    \textbf{Individual Mode}\\
    Each parallel iteration $\to$ one independent GPU thread.\\
    \vspace{0.5em}
    \emph{``Assign each student a task. They work alone.''}

    \column{0.48\textwidth}
    \textbf{Group Mode}\\
    Each parallel iteration $\to$ one GPU threadgroup.\\
    Threads inside the group coordinate.\\
    \vspace{0.5em}
    \emph{``Divide students into groups. Students in a group can collaborate.''}
  \end{columns}

  \vspace{1.5em}
  \begin{center}
  Both modes use the \emph{same} user-facing annotation:\\[0.5em]
  \mintinline{python}{#pragma parallel for}
  \end{center}
\end{frame}

\begin{frame}{Individual Mode: Concept}
  \begin{block}{Analogy}
    Assign each student an independent task. They work completely independently —
    no communication, no coordination. Each one just does their own thing.
  \end{block}

  \vspace{1em}
  On the GPU: each iteration of the \texttt{\#pragma parallel for} loop maps to
  \textbf{exactly one GPU thread}.

  \vspace{1em}
  Suitable for:
  \begin{itemize}
    \item Embarrassingly parallel workloads: Mandelbrot, ray marching, prime counting
    \item Workloads where iterations are fully independent
    \item At most a \emph{final} global reduction across all iterations
  \end{itemize}
\end{frame}

\begin{frame}[fragile]{Individual Mode: Mandelbrot Example}
  \begin{minted}{python}
  @appy.jit
  def mandelbrot(out, width, height, max_iter):
      #pragma parallel for
      for idx in range(width * height):
          x = idx % width
          y = idx // width
          c = complex(x / width * 3.5 - 2.5,
                      y / height * 2.0 - 1.0)
          z, i = 0+0j, 0
          while abs(z) <= 2.0 and i < max_iter:
              z = z * z + c
              i += 1
          out[y, x] = i
  \end{minted}
  Each pixel is computed by one independent thread.
  No communication between iterations.
\end{frame}

\begin{frame}{Group Mode: Concept}
  \begin{block}{Analogy}
    Divide students into groups and launch $X$ groups. Students \emph{within} a group
    can coordinate — they share a whiteboard, can communicate, and work together
    on their group's piece of the problem.
  \end{block}

  \vspace{1em}
  On the GPU: each iteration of the \texttt{\#pragma parallel for} loop maps to
  \textbf{one GPU threadgroup}. Threads inside a group coordinate via shared memory.

  \vspace{1em}
  Suitable for:
  \begin{itemize}
    \item Iterations whose body contains inner collective operations
    \item Dot products, slice reductions, scatter/gather over a dimension
    \item Kernels most naturally expressed with nested parallelism
  \end{itemize}
\end{frame}

\begin{frame}[fragile]{Group Mode: Symmetric Matrix Multiply (SYMM)}
  \begin{minted}{python}
  @appy.jit
  def kernel_appy(alpha, beta, C, A, B):
      C *= beta
      for i in range(C.shape[0]):
          #pragma parallel for          <- launch one group per i
          for j in range(C.shape[1]):
              C[:i, j] += alpha * B[i, j] * A[i, :i]  # collective write
              temp2[j]  = B[:i, j] @ A[i, :i]          # dot product
          C[i, :] += alpha * B[i, :] * A[i, i] + alpha * temp2
      return C
  \end{minted}

  \vspace{0.5em}
  Thread $j$ handles column $j$ — but within the group all threads cooperate on
  the slice operations and dot product via shared memory and threadgroup reductions.
\end{frame}

\begin{frame}[fragile]{Group Mode: TRMM (Triangular Matrix Multiply)}
  \begin{minted}{python}
  @appy.jit
  def kernel_appy(alpha, A, B):
      M, N = B.shape
      for i in range(M):
          #pragma parallel for
          for j in range(N):
              B[i, j] += A[i+1:, i] @ B[i+1:, j]  # dot product per (i,j)
      B *= alpha
      return B
  \end{minted}

  \vspace{0.5em}
  The \texttt{@} (dot product) inside the pfor body is a collective operation —
  the group of threads cooperates to reduce over the slice dimension.
\end{frame}

\begin{frame}[fragile]{Individual vs.\ Group Mode: Side by Side}
  \begin{center}
  \begin{tabular}{lll}
    \toprule
    & \textbf{Individual Mode} & \textbf{Group Mode} \\
    \midrule
    pfor iteration maps to & 1 GPU thread & 1 GPU threadgroup \\
    Thread coordination & None & Shared memory + barriers \\
    Inner parallelism & None & SIMD loops / reductions \\
    Suitable for & Independent tasks & Collective operations \\
    Examples & Mandelbrot, saxpy & TRMM, SYMM, GEMVER \\
    \midrule
    User annotation & \mintinline{python}{#pragma parallel for} & \mintinline{python}{#pragma parallel for} \\
    \bottomrule
  \end{tabular}
  \end{center}

  \vspace{1em}
  \begin{center}
  \textbf{The user writes the same annotation. The compiler infers the mode.}
  \end{center}
\end{frame}

% ---------------------------------------------------------------
\section{Compiler Design}

\begin{frame}[fragile]{How the Compiler Infers the Mode}
  The compiler detects \textbf{vector operations} inside the pfor loop body:
  \begin{itemize}
    \item Array slice assignments: \mintinline{python}{C[:i, j] += ...}
    \item Reductions: \mintinline{python}{np.sum(A[i+1:, i] * B[i+1:, j])}
    \item Dot products: \mintinline{python}{A[i+1:, i] @ B[i+1:, j]}
  \end{itemize}

  \vspace{1em}
  If any are found $\to$ \textbf{Group mode}: launch a threadgroup per iteration.\\
  Otherwise $\to$ \textbf{Individual mode}: launch one thread per iteration.

  \vspace{1em}
  The key compiler pass is \texttt{vector\_op\_to\_loop}, which lowers these
  vector operations into explicit inner loops with SIMD annotations.
\end{frame}

\begin{frame}{Compiler Pipeline Overview}
  \begin{center}
  \begin{tikzpicture}[
    box/.style={draw, rounded corners, minimum width=3.2cm, minimum height=0.65cm,
                font=\footnotesize, fill=mLightBrown!20},
    arr/.style={->, thick}
  ]
    \node[box] (src) {Python source + \texttt{\#pragma}};
    \node[box, below=0.45cm of src] (fe1) {rewrite\_aug\_assign};
    \node[box, below=0.45cm of fe1] (fe2) {attach\_pragma};
    \node[box, below=0.45cm of fe2] (fe3) {rewrite\_vector\_dot};
    \node[box, below=0.45cm of fe3] (fe4) {vector\_op\_to\_loop};
    \node[box, below=0.45cm of fe4] (fe5) {attach\_types};
    \node[box, below=0.45cm of fe5] (dev) {device code passes};
    \node[box, below=0.45cm of dev] (out) {Metal / Triton kernel};

    \draw[arr] (src) -- (fe1);
    \draw[arr] (fe1) -- (fe2);
    \draw[arr] (fe2) -- (fe3);
    \draw[arr] (fe3) -- (fe4);
    \draw[arr] (fe4) -- (fe5);
    \draw[arr] (fe5) -- (dev);
    \draw[arr] (dev) -- (out);

    \node[right=0.3cm of fe4, font=\scriptsize, text=mDarkTeal] {$\leftarrow$ mode inference happens here};
    \node[right=0.3cm of dev, font=\scriptsize, text=mDarkTeal] {$\leftarrow$ rewrite\_simd\_reductions,};
    \node[right=0.3cm of dev, font=\scriptsize, text=mDarkTeal, yshift=-0.35cm] {\phantom{$\leftarrow$} guard\_scalar\_mem\_writes};
  \end{tikzpicture}
  \end{center}
\end{frame}

\begin{frame}[fragile]{Key Pass: \texttt{vector\_op\_to\_loop}}
  Converts array-level operations into explicit scalar loops.
  \vspace{0.5em}

  \textbf{Input:}
  \begin{minted}{python}
  B[i, j] += A[i+1:, i] @ B[i+1:, j]
  \end{minted}

  \textbf{After rewrite\_aug\_assign + rewrite\_vector\_dot:}
  \begin{minted}{python}
  B[i, j] = B[i, j] + np.sum(A[i+1:, i] * B[i+1:, j])
  \end{minted}

  \textbf{After vector\_op\_to\_loop:}
  \begin{minted}{python}
  __reduce_sum_var = 0.0
  for __i0 in range(M - (i+1)):      # pragma simd, reduction
      __reduce_sum_var += A[__i0 + i+1, i] * B[__i0 + i+1, j]
  B[i, j] = B[i, j] + __reduce_sum_var
  \end{minted}

  The inner loop is marked \textbf{simd + reduction} $\to$ compiler generates threadgroup reduction.
\end{frame}

\begin{frame}[fragile]{Key Pass: \texttt{ExtractReductionSubexprs}}
  A new sub-pass inside \texttt{vector\_op\_to\_loop} that enables compound reduction expressions.

  \vspace{0.5em}
  \textbf{Before (required workaround):}
  \begin{minted}{python}
  s = np.sum(A[i+1:, i] * B[i+1:, j])
  B[i, j] = B[i, j] + s
  \end{minted}

  \textbf{After (now supported directly):}
  \begin{minted}{python}
  B[i, j] += A[i+1:, i] @ B[i+1:, j]
  \end{minted}

  The pass extracts the embedded reduction into a standalone assignment
  before the main lowering step, making the pipeline handle it naturally.
\end{frame}

\begin{frame}[fragile]{Key Pass: \texttt{guard\_scalar\_mem\_writes}}
  In group mode, many threads execute the pfor body — but scalar writes
  (e.g.\ \texttt{B[i,j] = ...}) must only happen \emph{once} per group.

  \vspace{0.5em}
  The pass automatically wraps scalar writes with a lane guard:

  \begin{minted}{c}
  // Generated Metal kernel (group mode)
  if (lane == 0) {
      B[i * N + j] = B[i * N + j] + __reduce_sum_var;
  }
  \end{minted}

  \vspace{0.5em}
  \textbf{Correctness guarantee:} only thread 0 in each group performs the
  scalar write, after the threadgroup reduction completes.
\end{frame}

\begin{frame}[fragile]{Key Pass: \texttt{rewrite\_simd\_reductions}}
  Converts the annotated reduction loop into a Metal threadgroup reduction:

  \begin{minted}{c}
  // Scalar loop body
  __reduce_sum_var += A[__i0 + i+1, i] * B[__i0 + i+1, j];
  // ...after loop:
  threadgroup float __tg_reduce_sum_var[SIMD_WIDTH];
  __threadgroup_reduce(__reduce_sum_var,
                       __tg_reduce_sum_var, 'sum');
  \end{minted}

  The tree-reduce implementation is emitted by \texttt{unparse\_to\_cpp},
  using Metal's \texttt{simdgroup\_sum} and \texttt{threadgroup} primitives.
\end{frame}

% ---------------------------------------------------------------
\section{Supported Workloads \& Examples}

\begin{frame}{NPBench Kernels Supported}
  \begin{center}
  \begin{tabular}{llll}
    \toprule
    \textbf{Kernel} & \textbf{Mode} & \textbf{Key operation} & \textbf{Annotation} \\
    \midrule
    SAXPY / vec-add & Individual & Pointwise & \texttt{\#pragma parallel for} \\
    GELU            & Individual & Pointwise & \texttt{\#pragma parallel for} \\
    Mandelbrot      & Individual & Irregular & \texttt{\#pragma parallel for} \\
    GESUMMV         & Group & \texttt{np.sum} & \texttt{\#pragma parallel for} \\
    GEMVER          & Group & \texttt{@} (dot) & \texttt{\#pragma parallel for} \\
    TRMM            & Group & \texttt{@} (dot) & \texttt{\#pragma parallel for} \\
    SYMM            & Group & \texttt{@} + scatter & \texttt{\#pragma parallel for} \\
    Gramschmidt     & Group & \texttt{@} (dot) & \texttt{\#pragma parallel for} \\
    SYRK / SYR2K    & Group & Slice accumulate & \texttt{\#pragma parallel for} \\
    Jacobi-1D/2D    & Group & Stencil & \texttt{\#pragma parallel for} \\
    \bottomrule
  \end{tabular}
  \end{center}
\end{frame}

\begin{frame}[fragile]{Writing APPy vs.\ NumPy Reference}
  \textbf{NumPy reference (TRMM):}
  \begin{minted}{python}
  for i in range(B.shape[0]):
      for j in range(B.shape[1]):
          B[i, j] += np.dot(A[i+1:, i], B[i+1:, j])
  B *= alpha
  \end{minted}

  \textbf{APPy kernel:}
  \begin{minted}{python}
  @appy.jit
  def kernel(alpha, A, B):
      for i in range(B.shape[0]):
          #pragma parallel for
          for j in range(B.shape[1]):
              B[i, j] += A[i+1:, i] @ B[i+1:, j]
      B *= alpha
  \end{minted}

  \vspace{0.5em}
  The APPy version is essentially the NumPy version with one added pragma.
\end{frame}

% ---------------------------------------------------------------
\section{Discussion \& Future Work}

\begin{frame}{What We've Shown So Far}
  \begin{itemize}
    \item A unified programming model that supports two GPU parallelism paradigms
          with a \textbf{single annotation}
    \item A compiler that automatically infers the mode and generates correct,
          efficient GPU code
    \item Support for compound reduction expressions (\texttt{+=\ np.sum(\ldots)},
          \texttt{@} inside expressions)
    \item Backends: \textbf{Metal} (macOS) and \textbf{Triton} (NVIDIA/Linux)
    \item Validated on a range of NPBench kernels
  \end{itemize}
\end{frame}

\begin{frame}{Limitations \& Open Questions}
  \begin{itemize}
    \item \textbf{1D array expressions only} inside pfor bodies — no 2D slice ops yet
    \item \textbf{No auto-tuning} for threadgroup sizes (currently fixed)
    \item \textbf{\texttt{\#pragma simd}} is an internal annotation that could be
          renamed \texttt{\#pragma group} to better reflect the semantics
    \item How to handle \textbf{memory allocation} for intermediate arrays inside pfor bodies?
    \item Formal correctness argument for the guard pass is informal today
  \end{itemize}
\end{frame}

\begin{frame}{Future Work}
  \begin{itemize}
    \item \textbf{Rename \texttt{\#pragma simd} $\to$ \texttt{\#pragma group}}:
          SIMD is too narrow — group captures that threads \emph{cooperate},
          not just execute in lockstep
    \item \textbf{2D and higher-dimensional array expressions} inside pfor loops
    \item \textbf{Common backend passes}: consolidate shared logic (e.g.\
          \texttt{ExtractReductionSubexprs}) rather than duplicating per backend
    \item \textbf{Performance evaluation} on NPBench against NumPy, CuPy, and Triton baselines
    \item \textbf{Auto-tuning} for threadgroup size and inner loop tile size
  \end{itemize}
\end{frame}

\begin{frame}[fragile]{Summary}
  \begin{block}{APPy's key contribution}
    A Python GPU programming model with a \textbf{single pragma} that supports
    two complementary parallelism paradigms — individual and group mode —
    without requiring users to think about threads, blocks, or shared memory.
  \end{block}

  \vspace{1em}
  \textbf{Individual mode:} ``Each student works alone on their task.''\\
  \textbf{Group mode:} ``Students form groups and collaborate.''\\[0.5em]
  Both expressed with \mintinline{python}{#pragma parallel for}.

  \vspace{1em}
  \begin{center}
    \large Questions \& Feedback Welcome
  \end{center}
\end{frame}

\end{document}
